\begin{frame}{Finale Änderungsanträge}
\tiny
\vspace{-0.45cm}
\renewcommand{\arraystretch}{0.9}
\setlength{\tabcolsep}{2pt}
\setlength{\abovetopsep}{0pt}
\setlength{\belowrulesep}{0pt}
\begin{tabularx}{\textwidth}{@{}>{\centering\arraybackslash}m{1.55cm} >{\RaggedRight\arraybackslash}m{\dimexpr\linewidth-1.55cm-2\tabcolsep\relax}@{}}
\toprule
\multicolumn{1}{m{1.55cm}}{\centering\textbf{Fraktion}} & \multicolumn{1}{m{\dimexpr\linewidth-1.55cm-2\tabcolsep\relax}}{\centering\textbf{Änderungsantrag}} \\
\midrule
\includegraphics[width=2cm,height=1cm,keepaspectratio]{Bilder/EVP.png} & b) Im Falle von zu wenigen Soldaten und Soldatinnen \textcolor{red}{\sout{kann}}\textcolor{green}{\uline{wird}} eine Bedarfspflicht eingeführt \textcolor{green}{\uline{werden. c) Im Falle von zu wenigen Soldaten und Soldatinnen nach der Bedarfspflicht wird eine allgemeine Wehrpflicht eingeführt, deren Regularien durch die Mitgliedsstaaten erstellt }}werden\textcolor{green}{\uline{. Eine Verpflichtung zum Dienst an der Waffe ist unzulässig}}. \\
\includegraphics[width=2cm,height=1cm,keepaspectratio]{Bilder/SD.png} & b) Im Falle von zu wenigen Soldaten und Soldatinnen \textcolor{red}{\sout{kann}}\textcolor{green}{\uline{wird}} eine Bedarfspflicht eingeführt\textcolor{red}{\sout{ werden}}. \\
\includegraphics[width=2cm,height=1cm,keepaspectratio]{Bilder/RE.png} & bereits existierenden Standards der North Atlantic Treaty Organization (NATO) \textcolor{red}{\sout{nicht widersprechen}}\textcolor{green}{\uline{widersprechen, um eigene Systeme entwickeln zu können und die Interessen der EU unabhängig umzusetzen}}. \\
\includegraphics[width=2cm,height=1cm,keepaspectratio]{Bilder/PfE.png} & Streitkräfte und Ausrüstung in eine gemeinsame Europäische Armee zu übertragen, mit dem langfristigen Ziel \textcolor{red}{\sout{der kompletten Auflösung aller nationaler Streitkräfte;}}\textcolor{green}{\uline{teilweise nationale Streitkräfte für eine EU-Armee zur Verfügung zu stellen;}} i) 2030 10 Prozent, ii) 2040 25 \textcolor{red}{\sout{Prozent, iii) 2050 50 Prozent, iv) 2070 100 }}Prozent seiner Streitkräfte und militärischen Güter in die Europäische Armee übertragen. \\
\includegraphics[width=2cm,height=1cm,keepaspectratio]{Bilder/Left.png} & 1. Jeder EU-Mitgliedstaat soll sich verpflichten, einen Teil seiner nationalen Streitkräfte und Ausrüstung in eine gemeinsame Europäische Armee zu übertragen\textcolor{red}{\sout{, mit dem langfristigen Ziel der kompletten Auflösung aller nationaler Streitkräfte;}}\textcolor{green}{\uline{.}} \\
\bottomrule
\end{tabularx}
\end{frame}

